\section{Week 2 - Pressure at a Point (Isotropy \& Scalarity)}

Consider a microscopic, wedge-shaped fluid element suspended in an air-water system below a free surface, subject to a gravitational acceleration field $\vec{g}$:

\begin{center}
\begin{tikzpicture}[scale=1.5, >=Stealth]
    % Coordinates for the triangular fluid element wedge
    \coordinate (O) at (0,0);
    \coordinate (X) at (2,0);
    \coordinate (Z) at (0,1.5);
    
    % Draw the wedge
    \draw[thick, fill=primary!5] (O) -- (X) -- (Z) -- cycle;
    
    % Dimensions label
    \draw[<->] (0,-0.2) -- (2,-0.2) node[midway, below] {$dy$};
    \draw[<->] (-0.2,0) -- (-0.2,1.5) node[midway, left] {$dz$};
    \draw (1, 0.75) node[above right] {$ds$};
    
    % Angle theta
    \draw (0.4,0) arc (0:36.87:0.4);
    \node at (0.55,0.15) {$\theta$};
    
    % Force/Pressure vectors acting on faces
    \draw[->, line width=1.2pt, accent] (-0.6, 0.75) -- (0, 0.75) node[midway, above] {$p_y$};
    \draw[->, line width=1.2pt, accent] (1, -0.6) -- (1, 0) node[midway, right] {$p_z$};
    
    % Normal force on inclined face
    \coordinate (M) at (1, 0.75);
    \draw[->, line width=1.2pt, accent] ($(M) + (0.48, 0.64)$) -- (M) node[midway, above right] {$p_s$};
    
    % Gravity vector
    \draw[->, line width=1.2pt, secondary] (0.7, 0.5) -- (0.7, 0.1) node[below] {$\vec{g}$};
\end{tikzpicture}
\end{center}

By applying a localized static force balance ($\Sigma \vec{F} = \vec{0}$) on the element while explicitly neglecting shearing forces ($P = \frac{F}{A}$):
\begin{equation}
    p_z \cdot dx dy - p_s \cdot ds dx \cos(\theta) - \rho g dV = \rho dV a_z
\end{equation}

In the vertical $z$-direction, this force balance maps directly to:
\begin{equation}
    p_z - p_s - \rho g \frac{dz}{2} = \rho a_z \frac{dz}{2}
\end{equation}

Taking the continuous limit as the geometric dimensions vanish ($dx, dy, dz \to 0$):
\begin{equation}
    p_z = p_s
\end{equation}

Performing an identical scale analysis along the horizontal $y$-direction (where the localized component of gravitational weight is identically zero):
\begin{equation}
    p_y = p_s
\end{equation}

\begin{conceptbox}[Isotropy of Hydrostatic Pressure]
Because $p_y = p_z = p_s$ under the limit $dV \to 0$, pressure is identically uniform in all directions at a single spatial point. Consequently, pressure behaves strictly as a \textbf{scalar field}.
\end{conceptbox}

\subsection{Governing Equation for the Pressure Field}

\subsubsection{Taylor Series Expansion on an Infinitesimal Fluid Volume}

To derive the continuous differential spatial distribution of pressure, we examine an isolated differential fluid cube of volume $dV = dx dy dz$ centered around a point $O$:

\begin{center}
\begin{tikzpicture}[scale=2, x={(-0.35cm,-0.2cm)}, y={(1cm,0cm)}, z={(0cm,1cm)}, >=Stealth]
    % Box parameters
    \def\dx{1.2}
    \def\dy{1.4}
    \def\dz{1.2}
    
    % Vertices of front face (y = 0)
    \coordinate (O) at (0,0,0);
    \coordinate (A) at (\dx,0,0);
    \coordinate (B) at (\dx,0,\dz);
    \coordinate (C) at (0,0,\dz);
    
    % Vertices of back face (y = \dy)
    \coordinate (O2) at (0,\dy,0);
    \coordinate (A2) at (\dx,\dy,0);
    \coordinate (B2) at (\dx,\dy,\dz);
    \coordinate (C2) at (0,\dy,\dz);
    
    % Draw dashed hidden lines
    \draw[dashed, color=gray] (O) -- (O2);
    \draw[dashed, color=gray] (O2) -- (A2);
    \draw[dashed, color=gray] (O2) -- (C2);
    
    % Draw solid visible lines
    \draw[thick] (O) -- (A) -- (B) -- (C) -- cycle;
    \draw[thick] (A) -- (A2) -- (B2) -- (B);
    \draw[thick] (C) -- (C2) -- (B2);
    \draw[thick] (A2) -- (B2) -- (C2);
    
    % Dimension annotations
    \draw[<->] (0,0,-0.15) -- (\dx,0,-0.15) node[midway, below] {$dx$};
    \draw[<->] (\dx+0.1,0,0) -- (\dx+0.1,0,\dz) node[midway, right] {$dz$};
    \draw[<->] (\dx+0.1,0,0) -- (\dx+0.1,\dy,0) node[midway, below right] {$dy$};
    
    % Pressure forces acting on Left and Right faces along y-axis
    % Center of left face
    \draw[->, line width=1.5pt, accent] (0.5*\dx, -0.4, 0.5*\dz) -- (0.5*\dx, 0, 0.5*\dz) node[left=4pt, pos=0.1] {$F_L$};
    
    % Center of right face
    \draw[<-, line width=1.5pt, accent] (0.5*\dx, \dy, 0.5*\dz) -- (0.5*\dx, \dy+0.4, 0.5*\dz) node[right=4pt, pos=0.9] {$F_R$};
    
    % Center gravity vector
    \draw[->, line width=1.2pt, secondary] (0.5*\dx, 0.5*\dy, 0.5*\dz) -- (0.5*\dx, 0.5*\dy, 0.1) node[below] {$\vec{w}$};
\end{tikzpicture}
\end{center}

Using a first-order multivariate Taylor-series expansion, we define the localized pressures acting on the left ($P_L$) and right ($P_R$) faces bounding the $y$-axis:
\begin{align}
    P_R &= P + \left.\frac{\partial P}{\partial y}\right|_0 \frac{dy}{2} + \mathcal{O}(dy^2) \\
    P_L &= P - \left.\frac{\partial P}{\partial y}\right|_0 \frac{dy}{2} + \mathcal{O}(dy^2)
\end{align}

Converting these localized scalar fields into their corresponding differential force vectors ($F = P \cdot A$) acting across the profile area $dA = dx dz$:
\begin{align}
    F_R &= P_R \, dx dz = \left(P + \left.\frac{\partial P}{\partial y}\right|_0 \frac{dy}{2}\right) dx dz \\
    F_L &= P_L \, dx dz = \left(P - \left.\frac{\partial P}{\partial y}\right|_0 \frac{dy}{2}\right) dx dz
\end{align}

\subsubsection{Multi-Dimensional Vector Formulation}

Applying Newton's second law along the $y$-axis profile while strictly neglecting viscous shearing stresses:
\begin{equation}
    -F_R + F_L = \rho \, dV a_y \implies -\left.\frac{\partial P}{\partial y}\right|_0 dx dy dz = \rho \, dx dy dz \, a_y
\end{equation}

Evaluating the limit as $dV \to 0$ yields the continuous differential equation for the acceleration component:
\begin{equation}
    -\left.\frac{\partial P}{\partial y}\right|_0 = \rho a_y
\end{equation}

Extending this identical structural logic symmetrically across the orthogonal $x$ and $y$ spatial directions:
\begin{equation}
    -\frac{\partial P}{\partial x} = \rho a_x, \quad -\frac{\partial P}{\partial y} = \rho a_y
\end{equation}

Combining these spatial derivatives via the vector gradient operator $\nabla P$:
\begin{equation}
    -\nabla P = -\frac{\partial P}{\partial x}\hat{\imath} - \frac{\partial P}{\partial y}\hat{\jmath} - \frac{\partial P}{\partial z}\hat{k}
\end{equation}

Incorporating the body force vector of fluid weight acting strictly down the vertical axis ($\vec{w} = -\rho \, dV g \hat{k}$), Newton's second law formulated per unit volume becomes:

\begin{eqbox}{General Hydrodynamic Equation of Motion}
\begin{equation}
    -\nabla P - \rho g \hat{k} = \rho \vec{a}
\end{equation}
\end{eqbox}

\subsection{Hydrostatics: Fluids at Rest}

\subsubsection{The Hydrostatic Differential Equation}

For a fluid configured completely at rest, the global acceleration vector is identically null ($\vec{a} = \vec{0}$). The general equation of motion reduces to:
\begin{equation}
    -\nabla P - \rho g \hat{k} = \vec{0} \implies \frac{\partial P}{\partial x} = 0, \quad \frac{\partial P}{\partial y} = 0, \quad \frac{\partial P}{\partial z} = -\rho g
\end{equation}

\begin{expertbox}[Density Field Dependencies]
This shows mathematically that for a fluid at rest, the pressure field depends strictly on the vertical coordinate ($P = P(z)$). The global behavior of this field is governed by the structural equation of state for fluid density $\rho = \rho(P, T)$.
\end{expertbox}

\subsubsection{Incompressible Fluid Integration}

For liquids or fluids modeling an constant density profile ($\rho = \text{cst}$), the differential equation integrates explicitly across an arbitrary vertical distance $\Delta z = z_2 - z_1$:
\begin{equation}
    \int_{P_1}^{P_2} dP = \int_{z_1}^{z_2} -\rho g \, dz \implies \Delta P = -\rho g \Delta z
\end{equation}

Defining $h = z_0 - z$ as the downward vertical depth measured relative to a baseline reference height $z_0$ at pressure $P_0$:
\begin{equation}
    P - P_0 = \rho g h
\end{equation}

\subsubsection{Pressure Scale Conventions}

\begin{defbox}{Absolute vs. Gauge Pressure}
Pressure measurements are categorized based on their underlying reference datum:
\begin{itemize}
    \item \textbf{Absolute Pressure ($P_{\text{abs}}$):} Measured relative to a perfect physical vacuum ($P_{\text{vacuum}} = 0$).
    \item \textbf{Gauge Pressure ($P_{\text{gauge}}$):} Measured relative to the local ambient atmospheric pressure ($P_{\text{atm}}$). 
\end{itemize}
The mathematical relation connecting the two domains is defined by:
\begin{equation}
    P_{\text{abs}} = P_{\text{gauge}} + P_{\text{atm}}
\end{equation}
\end{defbox}